\begin{table}[t]
\centering
\caption{The Employment Effect of Automatic Wage Indexation}
\label{tab:main}
\begin{tabular}{lccccc}
\hline\hline
 & (1) & (2) & (3) & (4) & (5) \\
 & Baseline & Excl. & Binary & Private & Lagged \\
 & & Hospitality & Treatment & Sector & Treatment \\
\hline
Cum.\ indexation & -1.0909** & -0.9754** & & -1.5373** & -0.4934 \\
 & (0.421) & (0.4244) & & (0.6463) & (0.3953) \\
Early $\times$ Post & & & -0.0399 & & \\
 & & & (0.0289) & & \\
\hline
Sector FE & Yes & Yes & Yes & Yes & Yes \\
Quarter FE & Yes & Yes & Yes & Yes & Yes \\
Observations & 579 & 547 & 579 & 483 & 560 \\
Clusters & 19 & 18 & 19 & 16 & 19 \\
\hline\hline
\end{tabular}
\begin{tablenotes}[flushleft]
\small
\item \textit{Notes:} Dependent variable is log quarterly employment (thousands). ``Cum.\ indexation'' is the cumulative mandatory wage increase applied to each sector through its joint committee regime. Column~(1): baseline TWFE with sector and quarter fixed effects. Column~(2): excludes hospitality (NACE~I), which had a strong COVID rebound. Column~(3): binary treatment comparing early-indexing sectors (pivot-triggered and quarterly) to annual-January sectors. Column~(4): excludes public sector (NACE~O, P, Q). Column~(5): one-quarter lagged treatment. Standard errors clustered by NACE section in parentheses. $^{***}p<0.01$, $^{**}p<0.05$, $^{*}p<0.1$.
\end{tablenotes}
\end{table}

